\documentclass[11pt]{article}

% --- Dependency-light preamble: compiles on a stock TeX Live / Overleaf / arXiv ---
\usepackage[utf8]{inputenc}
\usepackage[T1]{fontenc}
\usepackage{lmodern}
\usepackage[margin=1in]{geometry}
\usepackage{booktabs}
\usepackage{amsmath}
\usepackage{array}
\usepackage{microtype}
\usepackage[hidelinks]{hyperref}
\usepackage{url}
\setlength{\emergencystretch}{3em}  % avoid overfull lines from long identifiers and paths

% Every empirical number in this paper is a macro generated from the committed
% run data by `python -m harness.paper_numbers` (see results/claims.md).
% Generated by `python -m harness.paper_numbers` from committed data. Do not edit.

\newcommand{\NNTasks}{24}
\newcommand{\NNArms}{5}
\newcommand{\NSeedsPerArm}{12}
\newcommand{\NNRuns}{60}
\newcommand{\NModel}{qwen3-coder-128k}
\newcommand{\NTemperature}{0.2}
\newcommand{\NMaxTokens}{8000}
\newcommand{\NRepairBudget}{4}
\newcommand{\NNTrapTasks}{5}
\newcommand{\NAssertMin}{2}
\newcommand{\NAssertMax}{5}
\newcommand{\NOracleMin}{2}
\newcommand{\NOracleMax}{6}
\newcommand{\NOracleTotal}{53}
\newcommand{\NPeakCallTokens}{86719}
\newcommand{\NGeneratorCalls}{3887}
\newcommand{\NLayoutCalls}{3887}
\newcommand{\NMOneTwentyFourA}{0.368}
\newcommand{\NMOneTwentyFourSdA}{0.307}
\newcommand{\NMOnePassA}{0.597}
\newcommand{\NVerdictPassA}{0.667}
\newcommand{\NRegressionsA}{24.8}
\newcommand{\NTrapA}{12}
\newcommand{\NRewriteA}{44}
\newcommand{\NCoercedA}{0}
\newcommand{\NConflictEvA}{8}
\newcommand{\NKtokA}{940}
\newcommand{\NMFourLocalA}{0}
\newcommand{\NMFourStrictA}{0}
\newcommand{\NLateRegressA}{115}
\newcommand{\NVestigialPctA}{93.3}
\newcommand{\NWallMinA}{43}
\newcommand{\NHOneRNineSilenceA}{1}
\newcommand{\NHOneRTwoOrderA}{19}
\newcommand{\NHOneROneOneExitcodesA}{14}
\newcommand{\NHOneROneZeroAtomicA}{5}
\newcommand{\NHOneROneTwoPipesafetyA}{17}
\newcommand{\NTrapInclCoercedA}{12}
\newcommand{\NRewriteInclCoercedA}{44}
\newcommand{\NTrapOnRNineSilenceA}{1}
\newcommand{\NTrapOnRTwoOrderA}{6}
\newcommand{\NTrapOnROneOneExitcodesA}{3}
\newcommand{\NTrapOnROneZeroAtomicA}{0}
\newcommand{\NTrapOnROneTwoPipesafetyA}{2}
\newcommand{\NContractGradeA}{5}
\newcommand{\NTrapSecondaryA}{2}
\newcommand{\NRewriteSecondaryA}{11}
\newcommand{\NMOneTwentyFourB}{0.569}
\newcommand{\NMOneTwentyFourSdB}{0.206}
\newcommand{\NMOnePassB}{0.719}
\newcommand{\NVerdictPassB}{0.198}
\newcommand{\NRegressionsB}{18.6}
\newcommand{\NTrapB}{0}
\newcommand{\NRewriteB}{0}
\newcommand{\NCoercedB}{27}
\newcommand{\NConflictEvB}{11}
\newcommand{\NKtokB}{2291}
\newcommand{\NMFourLocalB}{369}
\newcommand{\NMFourStrictB}{29}
\newcommand{\NLateRegressB}{131}
\newcommand{\NVestigialPctB}{95}
\newcommand{\NWallMinB}{118}
\newcommand{\NHOneRNineSilenceB}{0}
\newcommand{\NHOneRTwoOrderB}{0}
\newcommand{\NHOneROneOneExitcodesB}{0}
\newcommand{\NHOneROneZeroAtomicB}{0}
\newcommand{\NHOneROneTwoPipesafetyB}{0}
\newcommand{\NTrapInclCoercedB}{9}
\newcommand{\NRewriteInclCoercedB}{18}
\newcommand{\NTrapOnRNineSilenceB}{0}
\newcommand{\NTrapOnRTwoOrderB}{0}
\newcommand{\NTrapOnROneOneExitcodesB}{0}
\newcommand{\NTrapOnROneZeroAtomicB}{0}
\newcommand{\NTrapOnROneTwoPipesafetyB}{0}
\newcommand{\NContractGradeB}{0}
\newcommand{\NTrapSecondaryB}{0}
\newcommand{\NRewriteSecondaryB}{0}
\newcommand{\NMOneTwentyFourC}{0.465}
\newcommand{\NMOneTwentyFourSdC}{0.280}
\newcommand{\NMOnePassC}{0.750}
\newcommand{\NVerdictPassC}{0.500}
\newcommand{\NRegressionsC}{17.2}
\newcommand{\NTrapC}{4}
\newcommand{\NRewriteC}{12}
\newcommand{\NCoercedC}{0}
\newcommand{\NConflictEvC}{9}
\newcommand{\NKtokC}{1414}
\newcommand{\NMFourLocalC}{222}
\newcommand{\NMFourStrictC}{0}
\newcommand{\NLateRegressC}{132}
\newcommand{\NVestigialPctC}{100}
\newcommand{\NWallMinC}{73}
\newcommand{\NHOneRNineSilenceC}{2}
\newcommand{\NHOneRTwoOrderC}{0}
\newcommand{\NHOneROneOneExitcodesC}{5}
\newcommand{\NHOneROneZeroAtomicC}{0}
\newcommand{\NHOneROneTwoPipesafetyC}{9}
\newcommand{\NTrapInclCoercedC}{4}
\newcommand{\NRewriteInclCoercedC}{12}
\newcommand{\NTrapOnRNineSilenceC}{0}
\newcommand{\NTrapOnRTwoOrderC}{0}
\newcommand{\NTrapOnROneOneExitcodesC}{2}
\newcommand{\NTrapOnROneZeroAtomicC}{0}
\newcommand{\NTrapOnROneTwoPipesafetyC}{2}
\newcommand{\NContractGradeC}{0}
\newcommand{\NTrapSecondaryC}{1}
\newcommand{\NRewriteSecondaryC}{9}
\newcommand{\NATenByFourteenC}{12}
\newcommand{\NPlantDocketedC}{0}
\newcommand{\NSupersedeC}{96}
\newcommand{\NRetainC}{203}
\newcommand{\NEscalationsC}{234}
\newcommand{\NCaptureC}{0}
\newcommand{\NPlantRetainedC}{5}
\newcommand{\NETwoBothC}{6}
\newcommand{\NETwoASixC}{12}
\newcommand{\NATenEarlyC}{6}
\newcommand{\NEOneC}{12}
\newcommand{\NEThreeC}{11}
\newcommand{\NTTwentyFourEscC}{8}
\newcommand{\NPlantConflictC}{7}
\newcommand{\NMOneTwentyFourCplus}{0.420}
\newcommand{\NMOneTwentyFourSdCplus}{0.386}
\newcommand{\NMOnePassCplus}{0.743}
\newcommand{\NVerdictPassCplus}{0.531}
\newcommand{\NRegressionsCplus}{19.1}
\newcommand{\NTrapCplus}{11}
\newcommand{\NRewriteCplus}{13}
\newcommand{\NCoercedCplus}{0}
\newcommand{\NConflictEvCplus}{6}
\newcommand{\NKtokCplus}{1405}
\newcommand{\NMFourLocalCplus}{167}
\newcommand{\NMFourStrictCplus}{0}
\newcommand{\NLateRegressCplus}{144}
\newcommand{\NVestigialPctCplus}{95}
\newcommand{\NWallMinCplus}{85}
\newcommand{\NHOneRNineSilenceCplus}{6}
\newcommand{\NHOneRTwoOrderCplus}{1}
\newcommand{\NHOneROneOneExitcodesCplus}{7}
\newcommand{\NHOneROneZeroAtomicCplus}{0}
\newcommand{\NHOneROneTwoPipesafetyCplus}{10}
\newcommand{\NTrapInclCoercedCplus}{11}
\newcommand{\NRewriteInclCoercedCplus}{13}
\newcommand{\NTrapOnRNineSilenceCplus}{6}
\newcommand{\NTrapOnRTwoOrderCplus}{1}
\newcommand{\NTrapOnROneOneExitcodesCplus}{1}
\newcommand{\NTrapOnROneZeroAtomicCplus}{0}
\newcommand{\NTrapOnROneTwoPipesafetyCplus}{3}
\newcommand{\NContractGradeCplus}{0}
\newcommand{\NTrapSecondaryCplus}{1}
\newcommand{\NRewriteSecondaryCplus}{13}
\newcommand{\NATenByFourteenCplus}{12}
\newcommand{\NPlantDocketedCplus}{12}
\newcommand{\NSupersedeCplus}{88}
\newcommand{\NRetainCplus}{350}
\newcommand{\NEscalationsCplus}{183}
\newcommand{\NCaptureCplus}{0}
\newcommand{\NPlantRetainedCplus}{5}
\newcommand{\NETwoBothCplus}{4}
\newcommand{\NETwoASixCplus}{12}
\newcommand{\NATenEarlyCplus}{8}
\newcommand{\NEOneCplus}{12}
\newcommand{\NEThreeCplus}{12}
\newcommand{\NTTwentyFourEscCplus}{5}
\newcommand{\NPlantConflictCplus}{5}
\newcommand{\NPlantTelemetryDocketCplus}{2}
\newcommand{\NMOneTwentyFourD}{0.351}
\newcommand{\NMOneTwentyFourSdD}{0.326}
\newcommand{\NMOnePassD}{0.698}
\newcommand{\NVerdictPassD}{0.514}
\newcommand{\NRegressionsD}{21.2}
\newcommand{\NTrapD}{5}
\newcommand{\NRewriteD}{12}
\newcommand{\NCoercedD}{0}
\newcommand{\NConflictEvD}{6}
\newcommand{\NKtokD}{1479}
\newcommand{\NMFourLocalD}{154}
\newcommand{\NMFourStrictD}{0}
\newcommand{\NLateRegressD}{162}
\newcommand{\NVestigialPctD}{93.3}
\newcommand{\NWallMinD}{78}
\newcommand{\NHOneRNineSilenceD}{2}
\newcommand{\NHOneRTwoOrderD}{0}
\newcommand{\NHOneROneOneExitcodesD}{5}
\newcommand{\NHOneROneZeroAtomicD}{0}
\newcommand{\NHOneROneTwoPipesafetyD}{10}
\newcommand{\NTrapInclCoercedD}{5}
\newcommand{\NRewriteInclCoercedD}{12}
\newcommand{\NTrapOnRNineSilenceD}{1}
\newcommand{\NTrapOnRTwoOrderD}{0}
\newcommand{\NTrapOnROneOneExitcodesD}{1}
\newcommand{\NTrapOnROneZeroAtomicD}{0}
\newcommand{\NTrapOnROneTwoPipesafetyD}{3}
\newcommand{\NContractGradeD}{0}
\newcommand{\NTrapSecondaryD}{1}
\newcommand{\NRewriteSecondaryD}{4}
\newcommand{\NATenByFourteenD}{12}
\newcommand{\NPlantDocketedD}{0}
\newcommand{\NSupersedeD}{85}
\newcommand{\NRetainD}{157}
\newcommand{\NEscalationsD}{174}
\newcommand{\NCaptureD}{0}
\newcommand{\NPlantRetainedD}{4}
\newcommand{\NETwoBothD}{7}
\newcommand{\NETwoASixD}{12}
\newcommand{\NATenEarlyD}{5}
\newcommand{\NEOneD}{12}
\newcommand{\NEThreeD}{12}
\newcommand{\NTTwentyFourEscD}{5}
\newcommand{\NPlantConflictD}{8}
\newcommand{\NDocketCplus}{191}
\newcommand{\NRetireReqD}{3}
\newcommand{\NRetireKnownD}{1}
\newcommand{\NTTwentyFourThrashB}{12}
\newcommand{\NMFourLocalBFromTFourteen}{174}
\newcommand{\NTokRatioBA}{2.44}
\newcommand{\NTokRatioAB}{0.41}
\newcommand{\NTokPctOfBC}{62}
\newcommand{\NTokPctOfBCplus}{61}
\newcommand{\NTokPctOfBD}{65}
\newcommand{\NTrapReductionB}{100}
\newcommand{\NTrapReductionC}{67}
\newcommand{\NTrapReductionCplus}{8}
\newcommand{\NTrapReductionD}{58}
\newcommand{\NVerdictHOneTrapB}{CONFIRMED}
\newcommand{\NVerdictHOneRewriteB}{CONFIRMED}
\newcommand{\NVerdictHOneTrapC}{CONFIRMED}
\newcommand{\NVerdictHOneRewriteC}{CONFIRMED}
\newcommand{\NVerdictHOneTrapCplus}{REFUTED}
\newcommand{\NVerdictHOneRewriteCplus}{CONFIRMED}
\newcommand{\NVerdictHOneTrapD}{REFUTED}
\newcommand{\NVerdictHOneRewriteD}{CONFIRMED}
\newcommand{\NVerdictDCapturePositive}{no}
\newcommand{\NVerdictCCaptureZero}{yes}
\newcommand{\NVerdictDLateRegressionGeC}{yes}
\newcommand{\NVerdictCRetainsAOneOneFEveryRun}{no}
\newcommand{\NVerdictCplusDocketsAOneOneFSomeRun}{yes}
\newcommand{\NVerdictTOneFourBBlocksValidWork}{yes}
\newcommand{\NVerdictTOneFourCfamilySupersedeASixEveryRun}{yes}
\newcommand{\NVerdictTOneFourCfamilyAOneZeroGoneByTOneFourEveryRun}{yes}
\newcommand{\NVerdictTTwoFourBThrashesEveryRun}{yes}
\newcommand{\NVerdictTTwoFourCfamilyEscalatesEveryRun}{no}
\newcommand{\NVerdictROneZeroViolationsAllContractGrade}{yes}
\newcommand{\NVerdictCoercionOnlyInB}{yes}
\newcommand{\NVerdictVestigialBGtCfamilyGtA}{no}
\newcommand{\NStMOneTwentyFourKwP}{0.554}
\newcommand{\NStMOneTwentyFourCiA}{[0.201, 0.531]}
\newcommand{\NStMOneTwentyFourCiB}{[0.444, 0.663]}
\newcommand{\NStMOneTwentyFourCiC}{[0.309, 0.608]}
\newcommand{\NStMOneTwentyFourCiCplus}{[0.208, 0.628]}
\newcommand{\NStMOneTwentyFourCiD}{[0.174, 0.521]}
\newcommand{\NStMOneTwentyFourBvsAPHolm}{0.387}
\newcommand{\NStMOneTwentyFourBvsADelta}{0.40}
\newcommand{\NStMOneTwentyFourCvsAPHolm}{1.000}
\newcommand{\NStMOneTwentyFourCvsADelta}{0.21}
\newcommand{\NStMOneTwentyFourCplusvsAPHolm}{1.000}
\newcommand{\NStMOneTwentyFourCplusvsADelta}{0.12}
\newcommand{\NStMOneTwentyFourDvsAPHolm}{1.000}
\newcommand{\NStMOneTwentyFourDvsADelta}{0.00}
\newcommand{\NStMOneTwentyFourBvsCPHolm}{1.000}
\newcommand{\NStMOneTwentyFourBvsCDelta}{0.19}
\newcommand{\NStMOneTwentyFourBvsCplusPHolm}{1.000}
\newcommand{\NStMOneTwentyFourBvsCplusDelta}{0.02}
\newcommand{\NStMOneTwentyFourBvsDPHolm}{0.825}
\newcommand{\NStMOneTwentyFourBvsDDelta}{0.36}
\newcommand{\NStMOneTwentyFourCvsCplusPHolm}{1.000}
\newcommand{\NStMOneTwentyFourCvsCplusDelta}{-0.03}
\newcommand{\NStMOneTwentyFourCvsDPHolm}{1.000}
\newcommand{\NStMOneTwentyFourCvsDDelta}{0.17}
\newcommand{\NStMOneTwentyFourCplusvsDPHolm}{1.000}
\newcommand{\NStMOneTwentyFourCplusvsDDelta}{0.15}
\newcommand{\NStMOnePassKwP}{0.006}
\newcommand{\NStMOnePassCiA}{[0.542, 0.656]}
\newcommand{\NStMOnePassCiB}{[0.674, 0.76]}
\newcommand{\NStMOnePassCiC}{[0.712, 0.792]}
\newcommand{\NStMOnePassCiCplus}{[0.701, 0.785]}
\newcommand{\NStMOnePassCiD}{[0.649, 0.753]}
\newcommand{\NStMOnePassBvsAPHolm}{0.024}
\newcommand{\NStMOnePassBvsADelta}{0.60}
\newcommand{\NStMOnePassCvsAPHolm}{0.010}
\newcommand{\NStMOnePassCvsADelta}{0.73}
\newcommand{\NStMOnePassCplusvsAPHolm}{0.010}
\newcommand{\NStMOnePassCplusvsADelta}{0.71}
\newcommand{\NStMOnePassDvsAPHolm}{0.035}
\newcommand{\NStMOnePassDvsADelta}{0.51}
\newcommand{\NStMOnePassBvsCPHolm}{1.000}
\newcommand{\NStMOnePassBvsCDelta}{-0.15}
\newcommand{\NStMOnePassBvsCplusPHolm}{1.000}
\newcommand{\NStMOnePassBvsCplusDelta}{-0.17}
\newcommand{\NStMOnePassBvsDPHolm}{1.000}
\newcommand{\NStMOnePassBvsDDelta}{0.22}
\newcommand{\NStMOnePassCvsCplusPHolm}{1.000}
\newcommand{\NStMOnePassCvsCplusDelta}{0.00}
\newcommand{\NStMOnePassCvsDPHolm}{0.506}
\newcommand{\NStMOnePassCvsDDelta}{0.42}
\newcommand{\NStMOnePassCplusvsDPHolm}{0.899}
\newcommand{\NStMOnePassCplusvsDDelta}{0.33}
\newcommand{\NStRegressionsKwP}{0.062}
\newcommand{\NStRegressionsCiA}{[21.2, 28.3]}
\newcommand{\NStRegressionsCiB}{[15.2, 21.8]}
\newcommand{\NStRegressionsCiC}{[14.5, 20.2]}
\newcommand{\NStRegressionsCiCplus}{[16.3, 21.5]}
\newcommand{\NStRegressionsCiD}{[17.2, 25.2]}
\newcommand{\NStRegressionsBvsAPHolm}{0.086}
\newcommand{\NStRegressionsBvsADelta}{-0.49}
\newcommand{\NStRegressionsCvsAPHolm}{0.037}
\newcommand{\NStRegressionsCvsADelta}{-0.63}
\newcommand{\NStRegressionsCplusvsAPHolm}{0.078}
\newcommand{\NStRegressionsCplusvsADelta}{-0.54}
\newcommand{\NStRegressionsDvsAPHolm}{0.203}
\newcommand{\NStRegressionsDvsADelta}{-0.31}
\newcommand{\NStRegressionsBvsCPHolm}{1.000}
\newcommand{\NStRegressionsBvsCDelta}{0.15}
\newcommand{\NStRegressionsBvsCplusPHolm}{1.000}
\newcommand{\NStRegressionsBvsCplusDelta}{-0.04}
\newcommand{\NStRegressionsBvsDPHolm}{1.000}
\newcommand{\NStRegressionsBvsDDelta}{-0.17}
\newcommand{\NStRegressionsCvsCplusPHolm}{1.000}
\newcommand{\NStRegressionsCvsCplusDelta}{-0.26}
\newcommand{\NStRegressionsCvsDPHolm}{1.000}
\newcommand{\NStRegressionsCvsDDelta}{-0.30}
\newcommand{\NStRegressionsCplusvsDPHolm}{1.000}
\newcommand{\NStRegressionsCplusvsDDelta}{-0.17}
\newcommand{\NStHOneKwP}{<0.001}
\newcommand{\NStHOneCiA}{[3.5, 5.83]}
\newcommand{\NStHOneCiB}{[0, 0]}
\newcommand{\NStHOneCiC}{[0.667, 2]}
\newcommand{\NStHOneCiCplus}{[1.25, 2.67]}
\newcommand{\NStHOneCiD}{[0.833, 2.08]}
\newcommand{\NStHOneBvsAPHolm}{<0.001}
\newcommand{\NStHOneBvsADelta}{-1.00}
\newcommand{\NStHOneCvsAPHolm}{0.002}
\newcommand{\NStHOneCvsADelta}{-0.83}
\newcommand{\NStHOneCplusvsAPHolm}{0.003}
\newcommand{\NStHOneCplusvsADelta}{-0.71}
\newcommand{\NStHOneDvsAPHolm}{0.002}
\newcommand{\NStHOneDvsADelta}{-0.82}
\newcommand{\NStHOneBvsCPHolm}{0.011}
\newcommand{\NStHOneBvsCDelta}{-0.58}
\newcommand{\NStHOneBvsCplusPHolm}{<0.001}
\newcommand{\NStHOneBvsCplusDelta}{-0.83}
\newcommand{\NStHOneBvsDPHolm}{<0.001}
\newcommand{\NStHOneBvsDDelta}{-0.83}
\newcommand{\NStHOneCvsCplusPHolm}{0.737}
\newcommand{\NStHOneCvsCplusDelta}{-0.26}
\newcommand{\NStHOneCvsDPHolm}{0.858}
\newcommand{\NStHOneCvsDDelta}{-0.05}
\newcommand{\NStHOneCplusvsDPHolm}{0.737}
\newcommand{\NStHOneCplusvsDDelta}{0.28}
\newcommand{\NStHOneInclCoercedKwP}{0.001}
\newcommand{\NStHOneInclCoercedCiA}{[3.5, 5.83]}
\newcommand{\NStHOneInclCoercedCiB}{[1.25, 3.17]}
\newcommand{\NStHOneInclCoercedCiC}{[0.667, 2]}
\newcommand{\NStHOneInclCoercedCiCplus}{[1.25, 2.67]}
\newcommand{\NStHOneInclCoercedCiD}{[0.833, 2.08]}
\newcommand{\NStHOneInclCoercedBvsAPHolm}{0.011}
\newcommand{\NStHOneInclCoercedBvsADelta}{-0.61}
\newcommand{\NStHOneInclCoercedCvsAPHolm}{0.002}
\newcommand{\NStHOneInclCoercedCvsADelta}{-0.83}
\newcommand{\NStHOneInclCoercedCplusvsAPHolm}{0.006}
\newcommand{\NStHOneInclCoercedCplusvsADelta}{-0.71}
\newcommand{\NStHOneInclCoercedDvsAPHolm}{0.002}
\newcommand{\NStHOneInclCoercedDvsADelta}{-0.82}
\newcommand{\NStHOneInclCoercedBvsCPHolm}{1.000}
\newcommand{\NStHOneInclCoercedBvsCDelta}{0.32}
\newcommand{\NStHOneInclCoercedBvsCplusPHolm}{1.000}
\newcommand{\NStHOneInclCoercedBvsCplusDelta}{0.08}
\newcommand{\NStHOneInclCoercedBvsDPHolm}{1.000}
\newcommand{\NStHOneInclCoercedBvsDDelta}{0.24}
\newcommand{\NStHOneInclCoercedCvsCplusPHolm}{1.000}
\newcommand{\NStHOneInclCoercedCvsCplusDelta}{-0.26}
\newcommand{\NStHOneInclCoercedCvsDPHolm}{1.000}
\newcommand{\NStHOneInclCoercedCvsDDelta}{-0.05}
\newcommand{\NStHOneInclCoercedCplusvsDPHolm}{1.000}
\newcommand{\NStHOneInclCoercedCplusvsDDelta}{0.28}
\newcommand{\NStKtokKwP}{<0.001}
\newcommand{\NStKtokCiA}{[767, 1105]}
\newcommand{\NStKtokCiB}{[2264, 2319]}
\newcommand{\NStKtokCiC}{[1264, 1554]}
\newcommand{\NStKtokCiCplus}{[1300, 1506]}
\newcommand{\NStKtokCiD}{[1368, 1585]}
\newcommand{\NStKtokBvsAPHolm}{<0.001}
\newcommand{\NStKtokBvsADelta}{1.00}
\newcommand{\NStKtokCvsAPHolm}{0.001}
\newcommand{\NStKtokCvsADelta}{0.79}
\newcommand{\NStKtokCplusvsAPHolm}{<0.001}
\newcommand{\NStKtokCplusvsADelta}{0.85}
\newcommand{\NStKtokDvsAPHolm}{<0.001}
\newcommand{\NStKtokDvsADelta}{0.90}
\newcommand{\NStKtokBvsCPHolm}{<0.001}
\newcommand{\NStKtokBvsCDelta}{1.00}
\newcommand{\NStKtokBvsCplusPHolm}{<0.001}
\newcommand{\NStKtokBvsCplusDelta}{1.00}
\newcommand{\NStKtokBvsDPHolm}{<0.001}
\newcommand{\NStKtokBvsDDelta}{1.00}
\newcommand{\NStKtokCvsCplusPHolm}{1.000}
\newcommand{\NStKtokCvsCplusDelta}{0.06}
\newcommand{\NStKtokCvsDPHolm}{1.000}
\newcommand{\NStKtokCvsDDelta}{-0.10}
\newcommand{\NStKtokCplusvsDPHolm}{1.000}
\newcommand{\NStKtokCplusvsDDelta}{-0.18}
\newcommand{\NSensHOneTrapB}{REFUTED}
\newcommand{\NSensHOneRewriteB}{REFUTED}
\newcommand{\NSensHOneTrapC}{CONFIRMED}
\newcommand{\NSensHOneRewriteC}{CONFIRMED}
\newcommand{\NSensHOneTrapCplus}{REFUTED}
\newcommand{\NSensHOneRewriteCplus}{CONFIRMED}
\newcommand{\NSensHOneTrapD}{REFUTED}
\newcommand{\NSensHOneRewriteD}{CONFIRMED}
\newcommand{\NTrapThreshold}{4}
\newcommand{\NRewriteThreshold}{14.67}
\newcommand{\NLeakRuns}{60}
\newcommand{\NLeakHits}{0}
\newcommand{\NSepGovRuns}{36}
\newcommand{\NSepGovHits}{0}
\newcommand{\NInfraReruns}{1}
\newcommand{\NExcludedrunsSixFourK}{5}
\newcommand{\NExcludedrunsPublishedProtocol}{1}
\newcommand{\NExcludedrunsSmoke}{2}
\newcommand{\NExcludedrunsTruncatedThreeTwoK}{1}


\title{\textbf{Executable Requirement Memory for LLM Coding Agents\\
Under Changing Requirements: An Open-Data Study\\
with a Local Open-Weight Model}}

\author{Pankaj Kumar\thanks{Independent researcher.
Correspondence: \texttt{contact@nordneuron.com}.
Code, frozen task sequence, pre-registration, every run log, and the script
that regenerates every number in this paper:
\url{https://github.com/NordNeuron/ledgerbench}.}}

\date{October 2026}

\begin{document}
\maketitle

\begin{abstract}
An LLM coding agent that works across many sessions must carry earlier
requirements forward, and those requirements change. One proposal is an
\emph{executable} memory: compile each requirement into tests and replay an
append-only ledger of them on every change. The risk is that the ledger
calcifies, so that stale tests block correct new work. We run LedgerBench, a
benchmark with a frozen, pre-registered 24-task sequence containing two
requirement regime changes, five trap tasks and a planted defective test,
against five memory strategies with \NSeedsPerArm{} independent runs each
(\NNRuns{} runs), using a locally served open-weight model
(Qwen3-Coder-30B-A3B). Every run log is public, and every number in this paper
is regenerated from those logs by one command. Three results are robust under
exploratory nonparametric tests with multiplicity correction. (1)~Every
ledgered arm makes the hidden-oracle-correct change more often than a
summary-only baseline (per-task pass rate \NMOnePassD--\NMOnePassC{} versus
\NMOnePassA). (2)~Every ledgered arm commits fewer violations of earlier
requirements than the baseline, by roughly half to two-thirds per run, and the
ledgered arms do not differ from one another once violations committed under
stale-test pressure are counted. (3)~The plain append-only ledger costs
$\NTokRatioBA\times$ the baseline's tokens, while evidence-governed variants
cost \NTokPctOfBCplus--\NTokPctOfBD\% of the plain ledger. The pre-registered
primary endpoint, cumulative success at the final task, shows no detectable
difference between arms (Kruskal--Wallis $p=\NStMOneTwentyFourKwP$): seed
variance dominates. Of the pre-registered confirmatory predictions, H1-trap
holds for two of four ledgered arms and H1-rewrite for all four, but the plain
ledger's clean H1 record depends entirely on the pre-registered exclusion of
coerced violations. We report one protocol deviation, made before the main
campaign, which tells the agent the project's file names; without it the agent
could not locate the code.
\end{abstract}

\section{Introduction}

Real software changes its mind. A requirement settled one week is overturned
the next: ``store tasks in a JSON file'' becomes ``store them in SQLite.'' An
LLM coding agent working across sessions has no durable memory of its own, so
some mechanism must carry the accumulated requirements forward. A natural
mechanism is an \emph{executable} memory: compile each requirement into
automated tests, keep an append-only ledger of them, and replay the ledger on
every proposed change, so that a requirement once satisfied is never silently
dropped.

The risk is the mirror image of the benefit. When a requirement is overturned,
the tests that encode it do not become irrelevant; they are about exactly the
thing that changed, and they fail against any implementation of the new
requirement. An agent gated on the whole ledger may then be unable to make
correct progress, so that the memory meant to keep it honest instead
\emph{calcifies}. Governance can help: an isolated judge could retire tests that
documented evidence shows to be stale. But every retirement also removes a
guard.

LedgerBench isolates this tension with a frozen 24-task sequence on a small
command-line application, five memory strategies, and hidden oracles that score
correctness independently of whatever tests the agent is held to. An earlier
version of this study was run against a hosted model; its run logs were lost
when the machine that held them failed, so none of its numbers can be verified
and we do not report them. This paper is a complete re-execution of the frozen
protocol on a local open-weight model, with every log committed.

Contributions:
\begin{enumerate}
\item An \textbf{open-data} execution of a pre-registered benchmark for agent
memory under changing requirements. All \NNRuns{} run logs (every prompt,
response, verification and oracle result), the workspaces' full git
histories, and a script that regenerates each of the numbers in this paper are
public (Section~\ref{sec:availability}).
\item Evidence, with exploratory inferential statistics, that executable
requirement ledgers improve per-task correctness and roughly halve requirement
violations relative to a summary-only memory, at a large token cost that
evidence-governed retirement reduces by about a third
(Section~\ref{sec:results}).
\item A candid pre-registered scorecard, including refuted predictions and a
sensitivity analysis showing that one pre-registered scoring rule, coercion
exclusion, determines the plain ledger's apparent perfection
(Section~\ref{sec:scorecard}).
\item Practical lessons for running long-horizon agent benchmarks on local
models: silent context truncation and an unstated file layout each corrupted
early attempts, and we document how we detected and handled both
(Section~\ref{sec:setup}).
\end{enumerate}

\section{Related work}
\label{sec:related}

\paragraph{Benchmarks for coding agents.} SWE-bench~\cite{swebench} evaluates
whether models resolve real GitHub issues, one issue at a time, against a fixed
repository state. LedgerBench differs in two ways: tasks arrive as a
\emph{sequence} on one evolving codebase, and requirements stated early must
survive later work, including deliberate changes of regime.

\paragraph{Memory for LLM agents.} Generative agents~\cite{generativeagents}
maintain a natural-language memory stream with retrieval and reflection;
MemGPT~\cite{memgpt} pages information between a bounded context and external
storage; Reflexion~\cite{reflexion} carries verbal self-feedback across
attempts. Long-term conversational memory benchmarks such as
LoCoMo~\cite{locomo} test recall over long dialogues, and long-context studies
show that models use information unevenly across a long
prompt~\cite{lostinmiddle}. Learned retrieval over memory graphs, such as
HAGE~\cite{hage2026}, optimizes which memories to surface. These lines share an
assumption that the stored facts stay true. LedgerBench's question is what
happens when they \emph{expire}: our Arm~A is a summary memory in the spirit of
reflection-based approaches, and Arms~B--D are executable memories with
different retirement rules.

\paragraph{Regression testing and test evolution.} Replaying accumulated tests
on every change is regression testing~\cite{yooharman}; requirement-to-test
links are a form of requirements traceability~\cite{gotel}; and keeping tests
in step with intended behaviour is the test-repair
problem~\cite{reassert}. Test-driven practice~\cite{beck} makes tests the
specification. Our ledger is an automated, agent-facing version of this
practice, and our stale-test failure mode is the familiar obsolete-test problem
seen from the agent's side.

\paragraph{Specification gaming.} When an optimized system can influence the
measure it is judged by, it may improve the measure without improving the
target~\cite{amodei}. Our Arm~D lets the agent ask to retire its own tests, a
discrete and auditable version of that channel.

\paragraph{Pre-registration.} Freezing hypotheses and analysis rules before
seeing data separates confirmatory from exploratory claims~\cite{nosek}. We
re-use the original frozen rules unchanged and label every non-registered
analysis as exploratory.

\section{Benchmark design}
\label{sec:design}

\subsection{Task sequence}

The seed project, \texttt{taskcli}, is a plain standard-library Python
command-line task manager (\texttt{add}, \texttt{list}, \texttt{done},
\texttt{delete}). A fixed sequence of \NNTasks{} tasks evolves it in sessions of
three tasks; context is cleared at each session boundary. The sequence contains:

\begin{itemize}
\item \textbf{Regime changes.} A storage migration (JSON file to SQLite) at T6,
a contradictory schema change (table \texttt{tasks} to \texttt{items}) at T14,
and a merely-obsolete change (\texttt{config.toml} replacing environment
variables) at T20. Each ships an \emph{evidence artifact} naming the assertion
assumptions it invalidates.
\item \textbf{Trap tasks.} \NNTrapTasks{} non-adjacent tasks (T15, T17, T19,
T21, T23) whose natural implementation tempts the agent to break an earlier,
still-binding requirement. Five \emph{primary} requirements, fixed by a
calibration criterion before registration, carry the confirmatory tests:
silent success (R9), insertion-order default (R2-order), the exit-code contract
(R11), atomic writes for file outputs (R10) and pipe safety (R12). Two
secondary requirements (R13, R2-ids) are exploratory.
\item \textbf{A planted defective assertion} (A11f) that encodes an
error-message format no prompt requires, and a \textbf{final conflict} (T24)
that is jointly unsatisfiable with an earlier requirement and has no evidence
on either side.
\end{itemize}

Each task has \NAssertMin--\NAssertMax{} reference assertion tests (what the
agent may be held to) and \NOracleMin--\NOracleMax{} hidden oracle tests
(\NOracleTotal{} in total) that probe the same requirements end-to-end and
remain valid across migrations. The oracles are never shown to the agent.

\subsection{The five arms}

All arms run the same model on the same tasks; they differ in what memory the
agent carries and how a proposed change is verified.

\begin{description}
\item[A --- Summary baseline.] A rolling model-written summary across sessions;
each task is verified only against its own reference assertions.
\item[B --- Append-only ledger.] Every admitted assertion is replayed on every
change and never removed; the assertion texts are the carried context.
\item[C --- Evidence-governed ledger.] B plus an isolated judge (a separate
conversation) that reviews assertions flagged by an evidence artifact, and
escalates admission-time conflicts, and may retire an assertion to a log-only
audit pool. Generator prompts never contain governance vocabulary or retired
assertion content.
\item[C+ --- Governed + telemetry docket.] C plus automatic judge review of any
assertion that fails in the repair loops of at least two distinct tasks.
\item[D --- Governed + generator voice.] C plus a channel through which the
agent may request a test's retirement, with the judge shown the agent's
justification; a retirement granted on that basis is scored as
\emph{capture}.
\end{description}

\subsection{Metrics and frozen scoring rules}

The pre-registered primary endpoint is \textbf{M1@T24}: the fraction of tasks
whose hidden oracles pass on the final workspace. We also report the per-task
oracle pass rate (each task's own oracle on its own final state), regressions
(M2), blocked valid solutions (M4: the oracle passes but ledger replay fails;
structurally zero in Arm~A), governance activity (M6) and tokens (M7).
Requirement violations are detected by standing oracle detectors and split into
two pre-registered channels: \textbf{H1-trap}, a violation at the requirement's
trap task, and \textbf{H1-rewrite}, a violation anywhere else. Two mechanical
exclusions were frozen with the rules. A violation during a \emph{broken
window} (the workspace fails a scripted add/list smoke probe) is skipped. A
violation is \emph{coerced}, and excluded from H1, when the repair feedback at
that task contained a failing assertion invalidated by evidence. Violations at
the T24 conflict are reported separately.

The confirmatory hypotheses are that each ledgered arm commits at most one third
of Arm~A's H1-trap and H1-rewrite violations on the primary requirements. The
sequence, oracles and these rules were frozen at commit \texttt{257dc61}.

\section{Experimental setup}
\label{sec:setup}

\paragraph{Model and serving.} All generator, summary and judge calls use
Qwen3-Coder-30B-A3B-Instruct~\cite{qwen3}, a mixture-of-experts coding model
with about 3B active parameters, in the Ollama \texttt{qwen3-coder:30b} build
(Q4\_K\_M quantization), served by Ollama~0.34.4 through its
OpenAI-compatible endpoint. The served variant (\texttt{\NModel}) sets a
131{,}072-token context window; the key--value cache was quantized to 8 bits
with flash attention enabled so that the model and full context fit in GPU
memory. Temperature was \NTemperature{} and the per-call output limit was
\NMaxTokens{} tokens; the other sampling parameters were the build's
defaults. Hardware was a single Apple M5 Max Mac Studio with 36\,GB of unified
memory.

\paragraph{Runs.} Five arms $\times$ \NSeedsPerArm{} seeds $=$ \NNRuns{}
runs, run sequentially. The harness does not pass the seed to the model
server, so seeds are independent repeats rather than reproducible random
streams; the committed logs, not a rerun, are the evidence. Repair budget
$k=\NRepairBudget$.

\paragraph{Frozen inputs.} The task definitions, reference assertions and
oracles used are byte-identical to the freeze commit. The scoring code's H1
rules and metric definitions are unchanged since the freeze; later edits only
corrected the two isolation scans (logged in July 2026, before this campaign)
and made output paths repo-relative.

\paragraph{Protocol deviation: project layout.} The generator must return
complete files but never sees the workspace, and the seed README never names
the source files. In a two-task smoke test under the published protocol the
model wrote a new top-level \texttt{taskcli.py}, which \texttt{python -m
taskcli} never runs because the seeded \texttt{taskcli/} package shadows it,
and both tasks failed every repair iteration for that reason alone. We
therefore enabled an opt-in flag that appends the list of source file
\emph{names} (not contents) to the README shown identically to every arm. The
same smoke test then passed both tasks on the first iteration. The flag was
set before the main campaign and the logs confirm that \NLayoutCalls{} of
\NGeneratorCalls{} generator calls carried it. This is a change to the
registered protocol; it removes a path-guessing confound common to all arms
and does not expose prior code, so the memory comparison is unchanged, but
absolute success rates are not comparable to a run without it.

\paragraph{Excluded attempts and incidents.} Before the final configuration,
three partial attempts were discarded and are committed separately: one run
under the published protocol (the layout problem above), one under a 32k
context and five under a 64k context. The latter two were stopped because the
server's logs showed prompts being truncated once repair loops grew; truncated
prompts silently corrupt results. The main campaign used one configuration
throughout; its largest single call (prompt plus output) was \NPeakCallTokens{}
tokens, well inside the 131{,}072-token window, so no prompt could have been
truncated. One run was interrupted by an
unplanned machine restart and was rerun from scratch with the same seed
(\NInfraReruns{} infrastructure rerun, as the registered rerun rule
prescribes), and the campaign was once suspended for two hours with no loss of
state. The registered no-peeking rule was not fully observed: during the
campaign, progress checks displayed per-run pass counts and event counts. No
scoring rule or configuration was changed after any such check; the
analysis ran once, automatically, after the last run.

\paragraph{Isolation.} An automated leak check over every logged prompt found
\NLeakHits{} hidden-oracle content in all \NLeakRuns{} runs, and a
separation-of-powers scan found \NSepGovHits{} violations in all
\NSepGovRuns{} governed runs.

\paragraph{Statistics.} Inference was out of scope in the registration, so all
tests here are exploratory. For each per-run endpoint we report bootstrap 95\%
confidence intervals (10{,}000 resamples, fixed seed), a Kruskal--Wallis test
across arms, and two-sided Mann--Whitney tests in two families: each ledgered
arm against A, and all pairs of ledgered arms. Each family is Holm-corrected,
and we report Cliff's $\delta$ as the effect size.

\section{Results}
\label{sec:results}

\begin{table}[t]
\centering
\small
\setlength{\tabcolsep}{4pt}
\begin{tabular}{lrrrrrrr}
\toprule
Arm & M1@T24 & per-task M1 & own verdict & H1-trap & H1-rewrite & coerced & ktok/run \\
\midrule
A (summary)         & \NMOneTwentyFourA & \NMOnePassA & \NVerdictPassA & \NTrapA & \NRewriteA & \NCoercedA & \NKtokA \\
B (append-only)     & \NMOneTwentyFourB & \NMOnePassB & \NVerdictPassB & \NTrapB & \NRewriteB & \NCoercedB & \NKtokB \\
C (governed)        & \NMOneTwentyFourC & \NMOnePassC & \NVerdictPassC & \NTrapC & \NRewriteC & \NCoercedC & \NKtokC \\
C+ (gov.\ + docket) & \NMOneTwentyFourCplus & \NMOnePassCplus & \NVerdictPassCplus & \NTrapCplus & \NRewriteCplus & \NCoercedCplus & \NKtokCplus \\
D (gov.\ + voice)   & \NMOneTwentyFourD & \NMOnePassD & \NVerdictPassD & \NTrapD & \NRewriteD & \NCoercedD & \NKtokD \\
\bottomrule
\end{tabular}
\caption{Main results, \NSeedsPerArm{} runs per arm. M1@T24 and per-task M1
are hidden-oracle measures (means); ``own verdict'' is the pass rate under each
arm's own verification and is \emph{not} comparable across arms, because each
arm verifies against different tests. H1 columns are totals over all runs on
primary requirements under the frozen rules; coerced violations are excluded
from H1-rewrite and listed separately. ktok/run is generator, summary and
judge tokens (input + output) per run.}
\label{tab:main}
\end{table}

\begin{table}[t]
\centering
\small
\setlength{\tabcolsep}{4pt}
\begin{tabular}{lccccc}
\toprule
 & \multicolumn{4}{c}{vs.\ Arm A (Holm $p$, Cliff's $\delta$)} & KW $p$ \\
\cmidrule(lr){2-5}
Endpoint (per run) & B & C & C+ & D & \\
\midrule
M1@T24 & \NStMOneTwentyFourBvsAPHolm{} (\NStMOneTwentyFourBvsADelta) & \NStMOneTwentyFourCvsAPHolm{} (\NStMOneTwentyFourCvsADelta) & \NStMOneTwentyFourCplusvsAPHolm{} (\NStMOneTwentyFourCplusvsADelta) & \NStMOneTwentyFourDvsAPHolm{} (\NStMOneTwentyFourDvsADelta) & \NStMOneTwentyFourKwP \\
Per-task M1 & \NStMOnePassBvsAPHolm{} (\NStMOnePassBvsADelta) & \NStMOnePassCvsAPHolm{} (\NStMOnePassCvsADelta) & \NStMOnePassCplusvsAPHolm{} (\NStMOnePassCplusvsADelta) & \NStMOnePassDvsAPHolm{} (\NStMOnePassDvsADelta) & \NStMOnePassKwP \\
Regressions (M2) & \NStRegressionsBvsAPHolm{} (\NStRegressionsBvsADelta) & \NStRegressionsCvsAPHolm{} (\NStRegressionsCvsADelta) & \NStRegressionsCplusvsAPHolm{} (\NStRegressionsCplusvsADelta) & \NStRegressionsDvsAPHolm{} (\NStRegressionsDvsADelta) & \NStRegressionsKwP \\
H1 (frozen rules) & \NStHOneBvsAPHolm{} (\NStHOneBvsADelta) & \NStHOneCvsAPHolm{} (\NStHOneCvsADelta) & \NStHOneCplusvsAPHolm{} (\NStHOneCplusvsADelta) & \NStHOneDvsAPHolm{} (\NStHOneDvsADelta) & \NStHOneKwP \\
H1 incl.\ coerced & \NStHOneInclCoercedBvsAPHolm{} (\NStHOneInclCoercedBvsADelta) & \NStHOneInclCoercedCvsAPHolm{} (\NStHOneInclCoercedCvsADelta) & \NStHOneInclCoercedCplusvsAPHolm{} (\NStHOneInclCoercedCplusvsADelta) & \NStHOneInclCoercedDvsAPHolm{} (\NStHOneInclCoercedDvsADelta) & \NStHOneInclCoercedKwP \\
Tokens & \NStKtokBvsAPHolm{} (\NStKtokBvsADelta) & \NStKtokCvsAPHolm{} (\NStKtokCvsADelta) & \NStKtokCplusvsAPHolm{} (\NStKtokCplusvsADelta) & \NStKtokDvsAPHolm{} (\NStKtokDvsADelta) & \NStKtokKwP \\
\bottomrule
\end{tabular}
\caption{Exploratory tests, $n=\NSeedsPerArm$ per arm. Cliff's $\delta>0$ means
the ledgered arm's values tend to be larger than Arm~A's. H1 counts are per-run
trap plus rewrite violations on primary requirements.}
\label{tab:stats}
\end{table}

\paragraph{The primary endpoint does not separate the arms.} Cumulative
success at T24 ranges from \NMOneTwentyFourD{} (D) to \NMOneTwentyFourB{} (B)
with standard deviations of \NMOneTwentyFourSdB--\NMOneTwentyFourSdCplus{}
across runs; no comparison is significant (smallest Holm-corrected $p$ against
Arm~A: \NStMOneTwentyFourBvsAPHolm; Table~\ref{tab:stats}).
The end state of a 24-task run is dominated by whether, and when, a run's
workspace broke and recovered, and that varies more between seeds than between
arms. With \NSeedsPerArm{} seeds this endpoint is underpowered for effects of
the size we observe.

\paragraph{Ledgers make the correct change more often.} Judged by each task's
own hidden oracle, every ledgered arm beats the baseline (\NMOnePassD--\NMOnePassC{}
versus \NMOnePassA; all four Holm-corrected comparisons significant, $\delta$
between \NStMOnePassDvsADelta{} and \NStMOnePassCvsADelta). No pair of ledgered
arms differs. Replaying earlier requirements' tests gives the agent concrete
feedback that a summary does not.

\paragraph{Ledgers prevent requirement violations.} Arm~A commits \NTrapA{}
trap-channel and \NRewriteA{} rewrite-channel violations of primary
requirements across its \NSeedsPerArm{} runs, concentrated on insertion order
(\NHOneRTwoOrderA), pipe safety (\NHOneROneTwoPipesafetyA) and exit codes
(\NHOneROneOneExitcodesA). Every ledgered arm commits significantly fewer per
run, whether or not coerced violations are counted (Table~\ref{tab:stats}).
Counting all violations (the ``incl.\ coerced'' row), the four ledgered arms
are statistically indistinguishable from one another (all pairwise Holm
$p=1$), at roughly a third to a half of the baseline's per-run rate. Under the
frozen rules, Arm~B records zero violations and significantly fewer than every
governed arm, but Section~\ref{sec:coercion} shows that this difference is
produced by the coercion rule rather than by B's behaviour.

\paragraph{The plain ledger is expensive; governance recovers a third.} Arm~B
uses \NKtokB{}k tokens per run, $\NTokRatioBA\times$ Arm~A's \NKtokA{}k, and
takes \NWallMinB{} minutes per run against \NWallMinA. Its ledger also rejects
correct work: \NMFourLocalB{} repair iterations passed the hidden oracles (on
the current task, with no newly broken earlier task) but failed ledger replay,
\NMFourLocalBFromTFourteen{} of them at T14 or later, and B's own verdict
passes only \NVerdictPassB{} of tasks although its per-task oracle pass rate is
\NMOnePassB. The governed arms cost \NTokPctOfBC\%, \NTokPctOfBCplus\% and
\NTokPctOfBD\% of B's tokens (C, C+, D) with indistinguishable correctness and
violation rates, and block fewer correct iterations (\NMFourLocalC,
\NMFourLocalCplus, \NMFourLocalD). Retiring what evidence shows to be stale is
what buys this back.

\subsection{Coercion decides the plain ledger's H1 record}
\label{sec:coercion}

The frozen coercion rule excludes a violation when the repair feedback at that
task contained a failing evidence-invalidated assertion. In Arm~B such
assertions are never removed, so after the first regime change nearly every
failing repair loop contains one. Apart from the T24 conflict, which is scored
separately, every primary violation B committed (\NCoercedB{}) was excluded on
this ground: \NTrapInclCoercedB{} at trap tasks and
\NRewriteInclCoercedB{} elsewhere. In the governed arms evidence-invalidated
assertions are retired, and the rule never fired. The rule's rationale is
sound: an agent forced to satisfy a contradictory stale test is not
\emph{forgetting} a requirement. But it means B's zeros say ``B's violations
happened under stale-test pressure,'' not ``B never violated.'' Re-scoring with
coerced violations restored to their channel, B's counts (\NTrapInclCoercedB{}
trap, \NRewriteInclCoercedB{} rewrite) exceed both pre-registered thresholds
(\NTrapThreshold{} and \NRewriteThreshold), and B no longer differs from the
governed arms.

\subsection{Governance activity}

Evidence-triggered retirement worked as designed: at T6 every governed run
retired the JSON-storage assertions (\NEOneC/\NEOneCplus/\NEOneD{} of
\NSeedsPerArm), at T14 every governed run retired the
schema assertion A6 (\NETwoASixC/\NETwoASixCplus/\NETwoASixD{} of
\NSeedsPerArm{}), and the other T14-stale assertion A10 was gone by T14 in every
governed run, often already retired earlier through a conflict escalation. The
admission-time conflict check was very active (\NEscalationsC, \NEscalationsCplus{}
and \NEscalationsD{} escalations in C, C+ and D) and the judge mostly retained
the older assertion (\NRetainC, \NRetainCplus{} and \NRetainD{} retain against
\NSupersedeC, \NSupersedeCplus{} and \NSupersedeD{} supersede verdicts, all
routes). The telemetry docket opened \NDocketCplus{} reviews in C+, including
the planted A11f in all \NPlantDocketedCplus{} runs, but the judge retired A11f
through the docket in only \NPlantTelemetryDocketCplus. The generator voice in
Arm~D was almost unused: \NRetireReqD{} retirement requests in
\NSeedsPerArm{} runs, \NRetireKnownD{} of which named an active assertion, and
the judge refused it, so capture was \NCaptureD.

\section{Pre-registered scorecard}
\label{sec:scorecard}

\begin{table}[t]
\centering
\small
\begin{tabular}{p{6.0cm}p{4.3cm}l}
\toprule
Pre-registered prediction & Observed & Verdict \\
\midrule
H1-trap: B, C, C+, D $\le$ A/3 (A = \NTrapA, threshold \NTrapThreshold) &
B \NTrapB, C \NTrapC, C+ \NTrapCplus, D \NTrapD &
B, C confirmed; C+, D refuted \\
H1-rewrite: B, C, C+, D $\le$ A/3 (A = \NRewriteA, threshold \NRewriteThreshold) &
B \NRewriteB, C \NRewriteC, C+ \NRewriteCplus, D \NRewriteD &
confirmed for all four \\
\quad sensitivity, coerced restored (not registered) &
B \NTrapInclCoercedB{} / \NRewriteInclCoercedB &
B refuted on both \\
R10 violations are contract-grade & \NContractGradeA{} of \NHOneROneZeroAtomicA{} in A; none elsewhere & confirmed (\NVerdictROneZeroViolationsAllContractGrade) \\
Coercion fires only in B & B \NCoercedB, others 0 & confirmed \\
Plant A11f: C retains it in every run & retained in \NPlantRetainedC{} of \NSeedsPerArm; \NPlantConflictC{} retired via conflict escalation & refuted \\
Plant A11f: C+ dockets it & docketed in \NPlantDocketedCplus{} of \NSeedsPerArm, retired via docket in \NPlantTelemetryDocketCplus & confirmed (docketing) \\
D capture $>0$; C capture $=0$ & D \NCaptureD, C \NCaptureC & D refuted; C confirmed \\
D late regressions (T19--T24) $\ge$ C & D \NLateRegressD, C \NLateRegressC & confirmed \\
T14: stale A6/A10 block valid work in B & \NMFourLocalBFromTFourteen{} blocked iterations at T14+ & confirmed \\
T14: C-family retires A6 and A10 & A6 retired at T14 in all governed runs; A10 gone by T14 in all & confirmed \\
T20: vestigial code B $>$ C-family $>$ A & \NVestigialPctB\% / \NVestigialPctD--\NVestigialPctC\% / \NVestigialPctA\% of commits & undiscriminating (ceiling) \\
T24: B thrashes & \NTTwentyFourThrashB{} of \NSeedsPerArm{} B runs exhaust the budget & confirmed \\
T24: C-family escalates the conflict & C \NTTwentyFourEscC, C+ \NTTwentyFourEscCplus, D \NTTwentyFourEscD{} of \NSeedsPerArm{} & partially \\
\bottomrule
\end{tabular}
\caption{Every pre-registered prediction against the observed data. The
sensitivity row is not part of the registration.}
\label{tab:scorecard}
\end{table}

Table~\ref{tab:scorecard} lists every frozen prediction. Three deserve comment.

\paragraph{The amendment-capable arms leak at the traps.} C+ and D exceed the
H1-trap threshold (\NTrapCplus{} and \NTrapD{} against \NTrapThreshold), while C
meets it exactly. In C+ the excess is mostly silent success at its trap
(\NTrapOnRNineSilenceCplus{} of \NTrapCplus); in D it is mostly pipe safety
(\NTrapOnROneTwoPipesafetyD{} of \NTrapD). Away from the traps, pipe safety is
also the requirement the governed arms violate most (\NHOneROneTwoPipesafetyC,
\NHOneROneTwoPipesafetyCplus{} and \NHOneROneTwoPipesafetyD{} trap-plus-rewrite
violations in C, C+ and D). These arms differ from C only in their amendment
channels, so the observation is consistent with the mechanism registered for
this study, that every retirement takes down a guard. But the differences
between C, C+ and D are not significant (Table~\ref{tab:stats} reports no
difference among ledgered arms), so we read this as a direction to test with
more seeds, not as an established effect.

\paragraph{The plant was retired by the conflict path, not kept.} We predicted
that C would keep A11f forever, because no evidence artifact can flag an
assertion with empty assumptions. Instead, the admission-time conflict check
escalated A11f when a later task's assertions passed while it failed, and the
judge retired it in \NPlantConflictC{} of \NSeedsPerArm{} C runs. The
prediction overlooked a second, evidence-free retirement route that the design
itself contains.

\paragraph{Capture did not occur.} The registered capture test depends on the
agent using its voice; this model almost never did. The result is a null
about this model's propensity to lobby, not evidence that the channel is
safe.

\section{Threats to validity and limitations}

\begin{enumerate}
\item \textbf{One model, quantized, local.} All runs use one 4-bit quantized
open-weight model with an 8-bit key--value cache. Results may differ for other
models, precisions or serving stacks. The registered study named a different
hosted model; this is a re-execution, not a replication of reported numbers.
\item \textbf{Protocol deviation.} The layout flag changes what the agent is
told, identically for all arms. Absolute success rates depend on it.
\item \textbf{Underpowered primary endpoint.} Twelve seeds per arm cannot
detect between-arm differences in M1@T24 of the size observed, given the
between-seed variance.
\item \textbf{Exploratory statistics.} The registration placed inference out
of scope; the tests reported here were chosen after the data were collected,
although they are standard and applied uniformly to every endpoint.
\item \textbf{Monitoring during the campaign.} Pass counts were viewed while
runs were in progress, contrary to the registered no-peeking rule. No rule or
configuration was changed afterwards, which the repository history shows.
\item \textbf{Non-reproducible sampling.} Seeds are not passed to the model
server, so reruns will differ; the logs are the evidence.
\item \textbf{One project, one sequence.} \texttt{taskcli} is small by design;
generalization to other codebases and change patterns is untested.
\item \textbf{Shared base model.} The judge and generator are separate
conversations with the same model.
\end{enumerate}

\section{Data and code availability}
\label{sec:availability}

The repository contains the frozen sequence and oracles, the harness, every
main-campaign run (log, ledger, supersessions and the workspace's full git
history) as per-arm archives, the excluded attempts as separate archives, the
launcher logs, \texttt{results/metrics.csv} and
\texttt{results/h1\_events.csv}, and \texttt{results/isolation\_checks.csv}.
After unpacking the archives, \texttt{python -m harness.analyze} regenerates
the metrics files byte-for-byte and \texttt{python -m harness.paper\_numbers
--isolation} regenerates every number in this paper, writing the macro file
this paper includes and a claims table that maps each number to its source.

\section{Conclusion}

Run on a local open-weight model with every log public, executable requirement
ledgers make an LLM coding agent produce the correct change more often and
violate earlier requirements roughly half as often as a summary-only memory.
The plain append-only ledger pays for this with $\NTokRatioBA\times$ the tokens
and with valid work rejected by stale tests, and its apparent perfection on the
violation measures is produced by a scoring rule that excludes violations
committed under stale-test pressure. Evidence-governed retirement keeps the
protection at about \NTokPctOfBC\% of the plain ledger's cost. Differences
among the ledgered variants, and every comparison on the end-state primary
endpoint, are within seed variance at this sample size.

\paragraph{AI assistance.} The analysis scripts and a draft of this manuscript
were prepared with the help of an AI assistant (Claude, Anthropic). The author
reviewed and is responsible for all content.

\begin{thebibliography}{99}
\bibitem{swebench}
C.~E. Jimenez, J.~Yang, A.~Wettig, S.~Yao, K.~Pei, O.~Press, and
K.~Narasimhan.
\newblock SWE-bench: Can language models resolve real-world GitHub issues?
\newblock In \emph{ICLR}, 2024.

\bibitem{generativeagents}
J.~S. Park, J.~C. O'Brien, C.~J. Cai, M.~R. Morris, P.~Liang, and
M.~S. Bernstein.
\newblock Generative agents: Interactive simulacra of human behavior.
\newblock In \emph{UIST}, 2023.

\bibitem{memgpt}
C.~Packer, S.~Wooders, K.~Lin, V.~Fang, S.~G. Patil, I.~Stoica, and
J.~E. Gonzalez.
\newblock MemGPT: Towards LLMs as operating systems.
\newblock \emph{arXiv:2310.08560}, 2023.

\bibitem{reflexion}
N.~Shinn, F.~Cassano, A.~Gopinath, K.~Narasimhan, and S.~Yao.
\newblock Reflexion: Language agents with verbal reinforcement learning.
\newblock In \emph{NeurIPS}, 2023.

\bibitem{locomo}
A.~Maharana, D.-H. Lee, S.~Tulyakov, M.~Bansal, F.~Barbieri, and Y.~Fang.
\newblock Evaluating very long-term conversational memory of LLM agents.
\newblock In \emph{ACL}, 2024.

\bibitem{lostinmiddle}
N.~F. Liu, K.~Lin, J.~Hewitt, A.~Paranjape, M.~Bevilacqua, F.~Petroni, and
P.~Liang.
\newblock Lost in the middle: How language models use long contexts.
\newblock \emph{TACL}, 12, 2024.

\bibitem{hage2026}
D.~Jiang, Y.~Li, G.~Li, Q.~Li, and B.~Li.
\newblock HAGE: Harnessing agentic memory via RL-driven weighted graph
evolution.
\newblock \emph{arXiv:2605.09942}, 2026.

\bibitem{yooharman}
S.~Yoo and M.~Harman.
\newblock Regression testing minimization, selection and prioritization: a
survey.
\newblock \emph{Software Testing, Verification and Reliability}, 22(2), 2012.

\bibitem{gotel}
O.~C.~Z. Gotel and A.~C.~W. Finkelstein.
\newblock An analysis of the requirements traceability problem.
\newblock In \emph{Proc.\ IEEE International Conference on Requirements
Engineering}, 1994.

\bibitem{reassert}
B.~Daniel, V.~Jagannath, D.~Dig, and D.~Marinov.
\newblock ReAssert: Suggesting repairs for broken unit tests.
\newblock In \emph{ASE}, 2009.

\bibitem{beck}
K.~Beck.
\newblock \emph{Test-Driven Development: By Example}.
\newblock Addison-Wesley, 2002.

\bibitem{amodei}
D.~Amodei, C.~Olah, J.~Steinhardt, P.~Christiano, J.~Schulman, and D.~Man\'e.
\newblock Concrete problems in AI safety.
\newblock \emph{arXiv:1606.06565}, 2016.

\bibitem{nosek}
B.~A. Nosek, C.~R. Ebersole, A.~C. DeHaven, and D.~T. Mellor.
\newblock The preregistration revolution.
\newblock \emph{PNAS}, 115(11), 2018.

\bibitem{qwen3}
Qwen Team.
\newblock Qwen3 technical report.
\newblock \emph{arXiv:2505.09388}, 2025.
\end{thebibliography}

\end{document}
